\documentclass[openany]{book}

\usepackage{amsmath, booktabs}
\usepackage{makecell, outlines, graphicx}
\usepackage[export]{adjustbox}
\usepackage[a4paper, top=.5in, bottom=.5in, left=.3in, right=.3in]{geometry}

\begin{document}

\chapter*{Topic 1}

\begin{outline}
	\1 $\lambda = \dfrac{v}{f} \implies v = \lambda \cdot f$
	\1 Power gain or loss in dB $= 10 \log \left( \dfrac{P_2}{P_1} \right)$
	\2 $P_2$ = Measured or calculated power
	\2 $P_1$ = Reference power
	\1 Voltage gain or loss in dB $= 10 \log \left( \dfrac{V_2}{V_1} \right)$
	\1 dBm is dB level using 1-mW reference, i.e., $P_1 = 1\ \text{mW}$
	\1 Signal to Noise Ratio (SNR) $= \dfrac{S}{N} = \mathrm{\dfrac{Signal Power}{Noise Power}} = \dfrac{P_S}{P_N}$
	\1 SNR $= 10 \log \dfrac{P_S}{P_N}$
	\1 Noise Figure, NF $= 10 \log \dfrac{S_i /N_i}{S_o /N_o} = 10 \log (\text{NR or SNR})$
	\1 Sensitivity, $S = -174 \text{dBm} + NF + 10 \log (BW) + (\text{desired } S/N)$
	\2 -174 dBm is the thermal noise power at room temperature (290 K) in a 1 Hz bandwidth (BW)
	\2 NF is of the reciever
	\2 $10 \log (BW)$, Change in noise power due to increasing  the bandwidth above 1 Hz
\end{outline}

\chapter*{Topic 2}

\begin{outline}
	\1 $m = \dfrac{E_m}{E_c}$
	\1 $V_{\max} = E_c + E_m$
	\1 $V_{\min} = E_c - E_m$
	\1 $m = \dfrac{V_{\max} - V_{\min}}{V_{\max} + V_{\min}}$
	\1 $M = m \times 100$
	\1 $E_c = \dfrac{1}{2} ( V_{ \max } + V_{ \min } )$
	\1 $E_m = \dfrac{1}{2} ( V_{ \max } - V_{ \min } )$
	\1 Peak change in amplitude of the modulated wave = 2 Em. This is so by how modulation works.
\end{outline}

\chapter*{Topic 3}

\begin{outline}
	\1 Amplitude Modulation = (Message Signal + Constant voltage) $\times$ Carrier Signal
	\2 Message Signal, $V_m(t) = E_m \sin (\omega_m t)$
	\2 Constant Voltage = $E_c$
	\2 Carrier Signal, $V_c(t) = \sin (\omega_c t)$
	\2 Amplitude Modulation,

	\begin{flalign*}
		V_{AM}(t) & = (V_m(t) + E_c) \times V_c(t)                                    \\
		          & = (E_m \sin (\omega_m t) + E_c) \times \sin (\omega_c t)          \\
		          & = E_m \sin (\omega_m t) \sin (\omega_c t) + E_c \sin (\omega_c t)
	\end{flalign*}

	We have, $\sin a \sin b = \dfrac{\cos(a - b) - \cos(a + b)}{2}$. Therefore,

	$$V_{AM}(t) = E_c \sin(2\pi f_c t) + E_m \dfrac{\cos(2\pi (f_c - f_m)t) - \cos(2\pi (f_c + f_m)t)}{2}$$

	\1 $P_c = \dfrac{E^2_c}{2R}$
	\1 $P_{USB} = P_{LSB} = P_c \cdot \dfrac{m^2}{4}$
	\1 $P_t = P_c + P_{USB} + P_{LSB}$
	\2 $P_t = P_c \cdot ( 1 + \dfrac{m^2}{2} )$, simplification of above equations
	\1 $m_t = \sqrt{ ( \sum m_n^2 ) }$
\end{outline}

\begin{outline}
	\1 $\Delta f = k_f \cdot V_{i(\max)}$
	\1 $f_{out} = f_c + \Delta f$
	\1 $m_p = k_p \cdot V_{i(\max)}$
	\1 Analog Modulation
	\2 Amplitude Modulation (DSB-FC)
	\3 $V_{AM(t)} = E_c \sin \omega_c t + \dfrac{m E_c}{2} \cos [ (\omega_c - \omega_i) t ] = \dfrac{m E_c}{2} \cos [ (\omega_c - \omega_i) t ]$
	\3 $P_{total} = \dfrac{E_c^2}{2R} \cdot \left( 1 + \dfrac{m^2}{2}\right)$
	\3 $BW = 2\cdot f_i$
	\3 FM
	\4 $V_{AM(t)} = E_c \sin ( \omega_c t + m_f \sin \omega_i t )$
	\4 $m_f = \dfrac{\Delta f}{f_i} = \dfrac{k_f \cdot V_{i(\max)}}{f_i}$
	\2 Angle Modulation
	\3 $P_{total} = $
\end{outline}

\begin{table}[!ht]
	\centering
	\renewcommand{\arraystretch}{2} % Increases row height for readability
	\begin{tabular}{|l|c|c|c|}
		\hline
		\textbf{Modulation Technique} & \textbf{Spectrum}                                                                                         & \textbf{Bandwidth (BW)}          & \textbf{Power}                         \\ \hline
		DSB-FC                        & \includegraphics[width=0.2\textwidth, valign=m, margin=0 5pt]{formulas-img/dsb-fc.png} \label{fig:dsb-fc} & $2 f_m$                          & $P_c \left( 1 + \frac{m^2}{2} \right)$ \\ \hline
		DSB-SC                        &
		\includegraphics[width=0.2\textwidth, valign=m, margin=0 5pt]{formulas-img/dsb-sc.png} \label{fig:dsb-sc}
		                              & $2 f_m$                                                                                                   & $2 \cdot \frac{(E_m / 2)^2}{2R}$                                          \\ \hline
		SSB-SC                        &
		\includegraphics[width=0.2\textwidth, valign=m, margin=0 5pt]{formulas-img/ssb-sc.png} \label{fig:ssb-sc}
		                              & $f_m$                                                                                                     & $\frac{(E_m / 2)^2}{2R}$                                                  \\ \hline
	\end{tabular}
	\caption{Comparison of Amplitude Modulation Techniques}
\end{table}



\chapter*{Topic 5}

\begin{outline}
	\1 $\Delta = \dfrac{A_{\max} - A_{\min} }{ \text{Number of quantization levels}}$
	\1 Bitrate (bps) = Sampling frequency $\times$ Number of bits per second
	\1 Quantization index, $i = \text{round}\left( \dfrac{x - x_{\min}}{\Delta} \right)$, where $x$ is a given analog value, and `round' means round down (remove the decimal part and leave the whole number alone)
	\1 Quantization value, $x_q = x_{\min} + i\Delta + \dfrac{\Delta}{2}$
	\1 Error, $e = x - x_q$
\end{outline}

\chapter*{Topic 7 M-ary Modulation}

\begin{outline}
	\1 Symbol (Baud) rate, $R_s = \dfrac{R_b}{\log_2 (M)}$
	\1 Fundamental frequency, $f_m = \dfrac{R_s}{2}$
	\1 Power for all M-ary Phase Shift Keying, $P = \dfrac{V^2}{2R}$
\end{outline}

\begin{table}[h]
	\centering
	\begin{tabular}{lc}
		\toprule
		\textbf{Modulation}        & \textbf{Formula}                   \\
		\midrule
		Binary \& Linear (PSK/QAM) & $BW  = R_s$                        \\
		\midrule
		Frequency (FSK)            & \makecell{$BW = 2(\Delta f + f_m)$ \\ $BW \propto M$ }   \\
		\bottomrule
	\end{tabular}
	\caption{BW}
\end{table}

\chapter*{Topic 8 Ber And Bandwidth Efficiency}

\begin{outline}
	\1 $B_{\eta} = \dfrac{\text{transmission bit rate (bps)}}{\text{minimum bandwidth (Hz)}}$
	\1 $\text{BER} \approx \dfrac{\text{number of bit errors}}{\text{total number of bits transmitted}}$
	\2 BER: Bit error rate, also called Probability of error
	\2 For example, if a system has a BER of $10^{-5}$, this means that on average we can expect one bit error in every $10^5$ bits transmitted
	\2 Probability of error is a function of:
	\3 The carrier-to-noise power ratio $C/N$ (or, more specifically, the average energy per bit-to-noise power density ratio $E_b/N_o$)
	\3 The number of possible encoding conditions used (M-ary)
	\1 $C_{(dBm)} = 10 \log \dfrac{C_{(watts)}}{0.001}$
	\2 C: average carrier power (combined power of the carrier and its associated sidebands)
	\1 $N = KTB$
	\2 $N$: thermal noise power (watts)
	\2 $K$: Boltzmann's constant ($1.38 \times 10^{-23} \text{ J/K}$)
	\2 $T$: temperature (K)
	\2 $B$: bandwidth (Hz)
	\2 $N_{(dBm)} = 10 \log \dfrac{N_{(watts)}}{0.001}$
	\1 $\dfrac{C}{N} = \dfrac{C}{KTB}$
	\2 $\dfrac{C}{N} (dB) = 10 \log \dfrac{C}{N} = C_{(dBm)} - N_{(dBm)}$
	\1 $E_b = C T_b= \dfrac{C}{f_b}$ (J/bit)
	\2 $E_b$: energy of a single bit (joules per bit)
	\2 $T_b$: time of a single bit (seconds)
	\2 $C$: carrier power (watts)
	\2 $E_{b(dBJ)} = 10 \log E_b = 10 \log \dfrac{C}{f_b} = 10 \log C - 10 \log f_b$
	\1 $N_o = \dfrac{N}{B} = \dfrac{KTB}{B} = KT$ (W/Hz)
	\2 $N_o$: noise power density (watts per hertz)
	\2 $N$: thermal noise power (watts)
	\2 $B$: bandwidth (Hertz)
	\2 $N_{o(dBm)} = \dfrac{N_{(dBm)}}{B_{(dBm)}} = 10 \log \dfrac{N_{(Watts)}}{0.001 \times B_{(Hertz)}} = 10 \log \dfrac{N}{0.001} - 10 \log B = N_{(dBm)} - 10 \log B$
	\2 $N_{o(dBm)} = 10 \log \dfrac{K}{0.001} + 10 \log T$
	\1 $\dfrac{B}{f_b}$: noise bandwidth-to-bit rate ratio
	\2 $B$: noise bandwidth (Hz)
	\2 $f_b$: bit rate (bps)
	\1 $\dfrac{E_b}{N_o} = \dfrac{C}{N} \times \dfrac{B}{f_b}$
	\2 $E_b/N_o$: energy per bit-to-noise power density ratio
	\2 $\dfrac{E_b}{N_0} \, (dB) = 10 \log \dfrac{C}{N} + 10 \log \dfrac{B}{f_b}$
	\2 The $E_b/N_0$ ratio is simply the product of the carrier-to-noise power ratio and the noise bandwidth-to-bit rate ratio
	\2 Mathematically: $\dfrac{E_b}{N_o} = \dfrac{C/f_b}{N/B} = \dfrac{CB}{Nf_b} = \dfrac{C}{N} \times \dfrac{B}{f_b}$
	\2 For binary modulation, the bandwidth equals the bit rate, so $E_b/N_0$ is equal to $C/N$
	\2 $E_b/N_o$ is used to compare two or more digital modulation systems that use different transmission rates (bit rates), modulation schemes (FSK, PSK, QAM), or encoding techniques (M-ary)
	\2 $E_b/N_o$ is the ratio of the energy of a single bit to the noise power present in 1 Hz of bandwidth
	\2 $E_b/N_o$ normalizes all multiphase modulation schemes to a common noise bandwidth, allowing for a simpler and more accurate comparison of their error performance
\end{outline}

Lecture 21:

\begin{outline}
	\1 Bit error performance for multiphase digital modulation systems is directly related to the distance between points on a signal state-space diagram
	\1 Threshold points (TP) for M-PSK:
	\2 $TP = \pm \dfrac{\pi}{M}$
	\2 $M$: number of signal states
	\2 Error occurs when noise shifts the received signal beyond a threshold into another decision region
	\1 For larger $M$ (both M-PSK and M-QAM):
	\2 Fixing $E_b/N_o$ $\Rightarrow$ $P_e$ increases
	\2 Fixing $P_e$ $\Rightarrow$ more power required
	\1 16-QAM vs 16-PSK (same bandwidth efficiency of 4 bit/s/Hz):
	\2 16-PSK needs $E_b/N_o$ of 16 dB for $P_e = 10^{-4}$
	\2 16-QAM needs $E_b/N_o$ of 12 dB for $P_e = 10^{-4}$ $\Rightarrow$ more power efficient
	\2 QPSK needs only 8.5 dB but bandwidth efficiency drops to 2 (twice the BW needed)
	\1 BW vs power/BER trade-off: higher-order modulation (larger $M$) $\Rightarrow$ better bandwidth efficiency but worse power efficiency (higher BER for same power)
\end{outline}
\end{document}
